\documentclass[11pt,a4paper]{article}
\newcommand{\intestazione}{Dimero di spin --- Schema dei passi sotto rumore}
% =====================================================================
%  Preambolo comune ai documenti sul dimero (serie "che cosa abbiamo fatto")
%  Incluso con % =====================================================================
%  Preambolo comune ai documenti sul dimero (serie "che cosa abbiamo fatto")
%  Incluso con % =====================================================================
%  Preambolo comune ai documenti sul dimero (serie "che cosa abbiamo fatto")
%  Incluso con \input{preambolo_dimero} subito dopo \documentclass.
% =====================================================================
\usepackage[T1]{fontenc}
\usepackage[utf8]{inputenc}
\usepackage[main=italian,provide=*]{babel}
\usepackage{amsmath,amssymb,mathtools}
\usepackage{graphicx}
\usepackage{booktabs}
\usepackage{colortbl}
\usepackage{xcolor}
\usepackage{geometry}
\usepackage{placeins}
\usepackage{caption}
\usepackage{enumitem}
\usepackage{fancyhdr}
\usepackage{needspace}
\usepackage{tikz}
\usetikzlibrary{arrows.meta,positioning,shapes.geometric,calc,fit,backgrounds}
\usepackage{hyperref}
\hypersetup{colorlinks=true,linkcolor=blu,citecolor=blu,urlcolor=blu}

\geometry{margin=2.5cm}

% --- colori coerenti con quelli delle figure (palette Okabe-Ito)
\definecolor{blu}{HTML}{0072B2}
\definecolor{arancio}{HTML}{E69F00}
\definecolor{verdes}{HTML}{009E73}
\definecolor{rossos}{HTML}{D55E00}
\definecolor{violas}{HTML}{CC79A7}
\definecolor{grigetto}{HTML}{E8E8E8}

\captionsetup{font=small,labelfont=bf,width=0.92\textwidth}

\pagestyle{fancy}
\fancyhf{}
\fancyhead[L]{\small\itshape\intestazione}
\fancyhead[R]{\small\thepage}
\renewcommand{\headrulewidth}{0.3pt}

% --- notazione
\newcommand{\ket}[1]{\lvert #1\rangle}
\newcommand{\bra}[1]{\langle #1 \rvert}
\newcommand{\media}[1]{\langle #1 \rangle}
\newcommand{\vs}[1]{\mathbf{s}_{#1}}

% --- riquadro "in una riga": il messaggio da portare a casa
\newcommand{\messaggio}[1]{%
  \begin{center}
  \begin{minipage}{0.93\textwidth}
  \setlength{\fboxsep}{9pt}%
  \colorbox{grigetto}{\begin{minipage}{\dimexpr\textwidth-2\fboxsep}
  \small\textbf{In una riga.}\enspace #1
  \end{minipage}}
  \end{minipage}
  \end{center}
}

% --- riquadro "come lo sappiamo": la verifica numerica dietro un'affermazione
\newcommand{\verifica}[1]{%
  \begin{center}
  \begin{minipage}{0.93\textwidth}
  \small\textcolor{verdes}{\rule{2.2em}{1.1pt}}\enspace
  \textcolor{verdes}{\textbf{Come lo sappiamo.}}\enspace #1
  \end{minipage}
  \end{center}
  \medskip
}
 subito dopo \documentclass.
% =====================================================================
\usepackage[T1]{fontenc}
\usepackage[utf8]{inputenc}
\usepackage[main=italian,provide=*]{babel}
\usepackage{amsmath,amssymb,mathtools}
\usepackage{graphicx}
\usepackage{booktabs}
\usepackage{colortbl}
\usepackage{xcolor}
\usepackage{geometry}
\usepackage{placeins}
\usepackage{caption}
\usepackage{enumitem}
\usepackage{fancyhdr}
\usepackage{needspace}
\usepackage{tikz}
\usetikzlibrary{arrows.meta,positioning,shapes.geometric,calc,fit,backgrounds}
\usepackage{hyperref}
\hypersetup{colorlinks=true,linkcolor=blu,citecolor=blu,urlcolor=blu}

\geometry{margin=2.5cm}

% --- colori coerenti con quelli delle figure (palette Okabe-Ito)
\definecolor{blu}{HTML}{0072B2}
\definecolor{arancio}{HTML}{E69F00}
\definecolor{verdes}{HTML}{009E73}
\definecolor{rossos}{HTML}{D55E00}
\definecolor{violas}{HTML}{CC79A7}
\definecolor{grigetto}{HTML}{E8E8E8}

\captionsetup{font=small,labelfont=bf,width=0.92\textwidth}

\pagestyle{fancy}
\fancyhf{}
\fancyhead[L]{\small\itshape\intestazione}
\fancyhead[R]{\small\thepage}
\renewcommand{\headrulewidth}{0.3pt}

% --- notazione
\newcommand{\ket}[1]{\lvert #1\rangle}
\newcommand{\bra}[1]{\langle #1 \rvert}
\newcommand{\media}[1]{\langle #1 \rangle}
\newcommand{\vs}[1]{\mathbf{s}_{#1}}

% --- riquadro "in una riga": il messaggio da portare a casa
\newcommand{\messaggio}[1]{%
  \begin{center}
  \begin{minipage}{0.93\textwidth}
  \setlength{\fboxsep}{9pt}%
  \colorbox{grigetto}{\begin{minipage}{\dimexpr\textwidth-2\fboxsep}
  \small\textbf{In una riga.}\enspace #1
  \end{minipage}}
  \end{minipage}
  \end{center}
}

% --- riquadro "come lo sappiamo": la verifica numerica dietro un'affermazione
\newcommand{\verifica}[1]{%
  \begin{center}
  \begin{minipage}{0.93\textwidth}
  \small\textcolor{verdes}{\rule{2.2em}{1.1pt}}\enspace
  \textcolor{verdes}{\textbf{Come lo sappiamo.}}\enspace #1
  \end{minipage}
  \end{center}
  \medskip
}
 subito dopo \documentclass.
% =====================================================================
\usepackage[T1]{fontenc}
\usepackage[utf8]{inputenc}
\usepackage[main=italian,provide=*]{babel}
\usepackage{amsmath,amssymb,mathtools}
\usepackage{graphicx}
\usepackage{booktabs}
\usepackage{colortbl}
\usepackage{xcolor}
\usepackage{geometry}
\usepackage{placeins}
\usepackage{caption}
\usepackage{enumitem}
\usepackage{fancyhdr}
\usepackage{needspace}
\usepackage{tikz}
\usetikzlibrary{arrows.meta,positioning,shapes.geometric,calc,fit,backgrounds}
\usepackage{hyperref}
\hypersetup{colorlinks=true,linkcolor=blu,citecolor=blu,urlcolor=blu}

\geometry{margin=2.5cm}

% --- colori coerenti con quelli delle figure (palette Okabe-Ito)
\definecolor{blu}{HTML}{0072B2}
\definecolor{arancio}{HTML}{E69F00}
\definecolor{verdes}{HTML}{009E73}
\definecolor{rossos}{HTML}{D55E00}
\definecolor{violas}{HTML}{CC79A7}
\definecolor{grigetto}{HTML}{E8E8E8}

\captionsetup{font=small,labelfont=bf,width=0.92\textwidth}

\pagestyle{fancy}
\fancyhf{}
\fancyhead[L]{\small\itshape\intestazione}
\fancyhead[R]{\small\thepage}
\renewcommand{\headrulewidth}{0.3pt}

% --- notazione
\newcommand{\ket}[1]{\lvert #1\rangle}
\newcommand{\bra}[1]{\langle #1 \rvert}
\newcommand{\media}[1]{\langle #1 \rangle}
\newcommand{\vs}[1]{\mathbf{s}_{#1}}

% --- riquadro "in una riga": il messaggio da portare a casa
\newcommand{\messaggio}[1]{%
  \begin{center}
  \begin{minipage}{0.93\textwidth}
  \setlength{\fboxsep}{9pt}%
  \colorbox{grigetto}{\begin{minipage}{\dimexpr\textwidth-2\fboxsep}
  \small\textbf{In una riga.}\enspace #1
  \end{minipage}}
  \end{minipage}
  \end{center}
}

% --- riquadro "come lo sappiamo": la verifica numerica dietro un'affermazione
\newcommand{\verifica}[1]{%
  \begin{center}
  \begin{minipage}{0.93\textwidth}
  \small\textcolor{verdes}{\rule{2.2em}{1.1pt}}\enspace
  \textcolor{verdes}{\textbf{Come lo sappiamo.}}\enspace #1
  \end{minipage}
  \end{center}
  \medskip
}

\usepackage{pdflscape}
\usepackage{array}

\title{\bfseries Schema riassuntivo --- Parte 2\\[2pt]
\large Punti di lavoro, risultati chiave e collegamenti sotto rumore}
\author{Samuele Schiavi \\ \small Tesi triennale in Fisica, Università di Parma
--- Relatore: Prof.\ Alessandro Chiesa}
\date{}

\begin{document}
\maketitle
\thispagestyle{fancy}

Questo documento non introduce fisica nuova: raccoglie in un unico posto,
per ciascuno dei sette documenti già scritti (cinque stadi più due
estensioni facoltative), il punto di lavoro usato, il risultato chiave e
il collegamento con quanto viene dopo --- pensato per essere consultato,
non letto in ordine. Riassume quello che c'è davvero nei Documenti 5--11,
non anticipa nulla.

\section*{Attenzione: cosa misura $N^*$, e cosa non misura}

Un chiarimento che vale per tutta la Parte 2, non solo per il documento in
cui è stato scoperto (Documento 8) --- va tenuto presente leggendo
qualunque $N^*$ da qui in avanti.

$N^*=\arg\max_N|C(N)|$ (o $F(N)$) trova il punto in cui il \textbf{segnale
è più misurabile sotto rumore} --- non quello in cui la stima del
correlatore fisico $C(t)$ è più \textbf{accurata}. Per la fedeltà $F(N)$
(Documento 7) questa distinzione non si pone: $F$ è per costruzione un
overlap con lo stato bersaglio, limitato fra $0$ e $1$, e un $F$ più alto
è sempre un'approssimazione migliore. Per il modulo di un correlatore
$|C(N)|$ (Documento 8) invece no: è stato verificato che a $N^*=2$,
$t=2$, $|C(N^*)|=0.5894$ --- ma il vero valore fisico, dalla somma
spettrale esatta, è $|C_\text{esatto}(t{=}2)|=0.3177$, quasi la metà.
L'errore di Trotter a $N$ piccolo non disperde il risultato attorno al
valore vero: lo \textbf{gonfia sistematicamente} (confermato a rumore
nullo: $0.6574$, ancora quasi il doppio del vero valore --- non è colpa
del rumore).

\messaggio{Per gli scopi di tutta la Parte 2 --- fissare un compromesso
Trotter/rumore, verificare se $N^*$ si sposta sotto altre condizioni --- è
il criterio giusto: serve un punto in cui il segnale emerga chiaramente
dal rumore, non un valore fisicamente accurato. Ma $|C(N^*)|$ non va mai
letto, in nessuno dei documenti che seguono, come ``il correlatore vale
questo'' --- andrebbe corretto per questo effetto, non preso alla
lettera.}

\section{Tabella d'insieme --- i cinque stadi}

\renewcommand{\arraystretch}{1.35}
\begin{landscape}
\thispagestyle{empty}
\begin{center}
\small
\setlength{\tabcolsep}{3.5pt}
\begin{tabular}{@{}>{\bfseries\raggedright}p{2.5cm} p{4.18cm} p{4.18cm} p{4.18cm} p{4.18cm} p{4.18cm}@{}}
\toprule
\rowcolor{grigetto}
 & \textbf{5. Il modello} & \textbf{6. Il VQE} & \textbf{7. La dinamica} & \textbf{8. I correlatori} & \textbf{9. Lo scan} \\
\midrule
Obiettivo &
Definire il rumore: due canali (depolarizzante 1q/2q, readout simmetrico),
calibrati su hardware reale &
Degradazione di un VQE ideale (Documento 2) eseguito su simulatore
rumoroso &
$N^*$: compromesso fra errore di Trotter e rumore accumulato, per la
fedeltà rispetto allo stato bersaglio fisico &
$N^*$ per il modulo di un correlatore dinamico, con una misura reale
sull'ancilla (readout genuinamente campionato) &
Come si sposta $N^*$ al variare dei parametri del canale di rumore \\
\addlinespace
Punto/i di lavoro &
$\varepsilon_{1q}=2.9\times10^{-4}$, $\varepsilon_{2q}=3.8\times10^{-3}$,
$p_\text{readout}=2.3\times10^{-2}$ (mediane \texttt{ibm\_torino}) &
``test 2'' ($b/J{=}0.35$, $D/J{=}0.80$) --- stesso punto singolo di
Parte 1 &
``test 2'', $t=2$ --- stesso punto, non quello della dinamica di Parte 1
($D/J=1$) &
``test 2'', $t=2$; correlatore $C_{11}^{yz}$, \textbf{lo stesso} già
scelto in Parte 1 (non uno nuovo) &
``test 2''; scan indipendente su $\varepsilon_{1q}$, $\varepsilon_{2q}$,
$p_\text{readout}$ \\
\addlinespace
Risultato chiave &
contrazione di Bloch piccola ($<1\%$) ma non nulla; identità
CNOT$\equiv$CZ verificata a precisione macchina &
$1-F\sim1.5\%$ a rumore di riferimento; costo di preparazione $2$ CNOT
(il più piccolo di tutta la pipeline) &
$N^*=8$, $F(N^*)=0.817725$; minimo locale a $N=3$ ($F\approx0.40$) ---
Trotter puro, non il rumore &
$N^*=2$, $|C(N^*)|=0.5894$ --- che sovrastima il vero
$|C_\text{esatto}(t{=}2)|=0.3177$ di $\sim1.9\times$ (vedi riquadro sopra)
&
$N^*$ monotono non-crescente in $\varepsilon_{1q},\varepsilon_{2q}$;
invariante rispetto a $p_\text{readout}$ simmetrico (dimostrato
\emph{e} verificato) \\
\addlinespace
Collegamento &
fornisce il \texttt{NoiseModel} riusato, senza modifiche, in tutta la
Parte 2 &
fissa la scala di degradazione ``di base'' (sola preparazione) per il
passo successivo &
introduce $N^*$ e la regola della transpilazione a blocchi, riusata tale
e quale in Documento 8 &
fornisce il correlatore e l'$N^*$ di riferimento per i Documenti 9, 10, 11
&
apre le due estensioni facoltative (Documenti 10, 11) \\
\bottomrule
\end{tabular}
\end{center}
\end{landscape}

\clearpage

\section{Tabella d'insieme --- le due estensioni}

\renewcommand{\arraystretch}{1.35}
\begin{center}
\small
\setlength{\tabcolsep}{5pt}
\begin{tabular}{@{}>{\bfseries\raggedright}p{2.6cm} p{5.6cm} p{5.6cm}@{}}
\toprule
\rowcolor{grigetto}
 & \textbf{10. VQE noise-aware} & \textbf{11. Readout asimmetrico} \\
\midrule
Obiettivo &
Rioptimizzare l'ansatz vedendo il rumore, invece di riusare i parametri
ideali --- cambia qualcosa a valle? &
Il test opzionale del relatore: cosa succede a $N^*$ se $p_{01}\neq p_{10}$? \\
\addlinespace
Punto/i di lavoro &
``test 2'' \emph{più} un secondo punto ($b/J{=}{-}0.18$, $D/J{=}1$, DM
$5\times$ più forte), per generalità &
``test 2''; split illustrativo $p_{10}{:}p_{01}=3\times$ alla stessa
media --- nessun valore reale citabile trovato per \texttt{ibm\_torino} \\
\addlinespace
Risultato chiave &
nessun cambiamento misurabile ($\Delta E\sim10^{-7}$--$10^{-8}$); $N^*$
invariato ovunque testato, inclusa l'intera griglia del Documento 9
($25/25$ punti) &
$N^*$ invariato anche nel caso asimmetrico, verificato fino a un rapporto
$50\times$ --- \textbf{non} garantito dalla teoria (a differenza del caso
simmetrico), quindi un fatto verificato, non assunto \\
\addlinespace
Collegamento &
chiude con un NO la domanda ``vale la pena rioptimizzare sotto rumore?'' &
chiude la Parte 2 sul dimero; il passo successivo del progetto è la
stessa pipeline sul sistema a tre spin (anello) \\
\bottomrule
\end{tabular}
\end{center}

\section{Il modello di rumore}

Due canali soli, non uno stack generico: depolarizzante (separato a 1 e 2
qubit) e readout simmetrico, calibrati sulle mediane pubblicate per
\texttt{ibm\_torino} (Stenger \emph{et al.}, arXiv:2504.15187). $T_1/T_2$
esclusi per decisione esplicita del relatore, non per omissione. Un punto
tecnico verificato, non assunto: il gate nativo a 2 qubit di
\texttt{ibm\_torino} è il CZ, mentre il circuito di questa tesi usa CNOT
--- i due sono legati da un'identità circuitale esatta
($\text{CNOT}=(\mathbb1\otimes H)\cdot\text{CZ}\cdot(\mathbb1\otimes H)$),
verificata a precisione macchina, quindi riusare $\varepsilon_{2q}$ del CZ
per il CNOT non è un'approssimazione arbitraria.

\messaggio{Il modello è deliberatamente minimale: due canali, quattro
numeri, calibrati su un chip reale. Questa scelta ha una conseguenza
diretta su tutto quello che segue --- ``rumoroso'', in questa parte della
tesi, significa esattamente questi due canali a questi valori (o a quelli
esplicitamente scansionati in Documento 9), non una replica generica di
``un computer quantistico rumoroso''.}

\section{Il VQE sotto rumore}

Il documento più semplice della serie: stesso ansatz, stessi parametri
$\theta^*$ di Parte 1, nessuna rioptimizzazione --- misura quanto degrada
un risultato \emph{ideale} se eseguito su hardware rumoroso, non quanto
bene converga un VQE che veda il rumore durante l'ottimizzazione (quella
domanda è rimandata, per intero, al Documento 10). Il costo di
preparazione, $2$ CNOT, è il più piccolo di tutta la pipeline: lascia poco
spazio al rumore di accumularsi, ed è per questo che la perdita di
fedeltà misurata qui ($\sim1.5\%$) è la scala più piccola di tutta la
Parte 2 --- gli stadi successivi, con circuiti più lunghi, partono da
questa scala e la degradano ulteriormente.

\section{La dinamica sotto rumore}

Primo punto della pipeline in cui il numero di passi di Trotter $N$
smette di essere ``più è meglio'': con il rumore acceso, ogni passo in
più costa $3$ CNOT --- il compromesso fra errore di Trotter (decrescente
in $N$) e rumore accumulato (crescente in $N$) fissa un $N^*$ finito
($N^*=8$), non $N\to\infty$ come nel sistema chiuso. Un dettaglio tecnico
riusato ovunque da qui in avanti: il singolo passo va transpilato una
sola volta al suo costo minimo e composto $N$ volte --- transpilare
l'intero circuito assemblato fonde gli $N$ passi in un blocco a costo
costante, cancellando silenziosamente la dipendenza da $N$ che è
l'oggetto dello studio.

\messaggio{Il minimo locale a $N=3$ (non un semplice calo, un vero
minimo) è riprodotto identico a rumore nullo e a rumore acceso: è un
effetto di Trotter puro, non un artefatto del modello di rumore --- lo
stesso tipo di comportamento non perturbativo ricomparirà, in forma più
marcata, nel correlatore del Documento 8.}

\section{I correlatori sotto rumore}

Qui il rumore di lettura entra per la prima volta in gioco: a differenza
di energia e fedeltà (lette direttamente dall'operatore densità, nessuna
misura), il correlatore passa da una vera istruzione di misura
sull'ancilla. Il correlatore usato non è una nuova scelta indipendente:
è \textbf{lo stesso} $C_{11}^{yz}$ già scelto e motivato in Parte 1
(Documento 4), riusato qui per continuità diretta --- e lo scan
sistematico sulle 36 combinazioni (rifatto sotto rumore, verifica non
data per scontata) conferma che è un caso rappresentativo, non un
estremo: sta nel secondo gruppo più popoloso ($N^*=2$, $12$ combinazioni
su $36$).

\begin{center}
\small
\renewcommand{\arraystretch}{1.35}
\begin{tabular}{@{}p{3.4cm}p{3.0cm}p{7.0cm}@{}}
\toprule
\rowcolor{grigetto}
\textbf{Correlatore} & \textbf{$a_2/a_1$, $N^*$} & \textbf{Dove si colloca} \\
\midrule
$C_{11}^{yz}$ \textit{(usato qui)} &
$0.71$, $N^*=2$ &
Stesso di Parte 1; secondo gruppo più popoloso su cinque --- rappresentativo,
non un estremo \\
$C_{21}^{xx}$ \textit{(usato in una versione precedente)} &
$0.48$, $N^*=5$ &
L'estremo più fragile di tutte le 36 combinazioni --- $N^*$ più alto,
$|C(N^*)|$ più basso in assoluto \\
Gruppo più ricco in assoluto &
$0.96$, $N^*=1$ &
Non usato: il massimo teorico di ricchezza spettrale non implica il
comportamento più interessante sotto rumore \\
\bottomrule
\end{tabular}
\end{center}

\messaggio{Il meccanismo dietro la posizione di $N^*$, verificato su
tutte e 36 le combinazioni: non è la ricchezza spettrale $a_2/a_1$ a
deciderla direttamente, è se la dinamica di Trotter ha o no un guadagno
non perturbativo disponibile appena oltre $N=1$. Le combinazioni con
$N^*=1$ sono già vicine al proprio plateau asintotico fin da $N=1$; le
altre hanno un guadagno sostanziale da inseguire (fino a $+1229\%$ nel
caso più estremo) --- vale la pena pagare qualche CNOT in più.}

\section{Lo scan sui parametri di rumore}

L'unico documento dei cinque senza un equivalente diretto in Parte 1:
un punto di calibrazione solo (Documenti 6--8) non dice come si comporti
il risultato altrove. Una previsione scritta \emph{prima} di guardare i
dati --- $N^*$ deve dipendere dal rumore di gate, non da quello di
lettura simmetrico, perché quest'ultimo è un fattore puramente
moltiplicativo e non può mai spostare un argmax --- è stata confermata
su entrambe le metriche (fedeltà di Trotter, modulo del correlatore),
non solo argomentata per plausibilità.

\messaggio{Un limite dichiarato esplicitamente: la verifica di
invarianza dal readout usa la stessa formula/lo stesso meccanismo
moltiplicativo che la garantisce teoricamente --- per il caso simmetrico
questo non è un problema (l'argomento matematico è comunque valido), ma
è la ragione precisa per cui il caso \emph{asimmetrico} (Documento 11),
dove quella garanzia matematica non c'è più, va verificato numericamente
da zero, non dedotto per analogia da qui.}

\section{Le due estensioni}

Due documenti facoltativi, tenuti separati perché non condividono
argomento fisico --- uno riguarda una strategia di ottimizzazione,
l'altro un'assunzione diversa nel modello di rumore. Condividono solo
l'esito (``non cambia $N^*$''), un filo troppo debole per giustificare di
comprimerli insieme.

\textbf{VQE noise-aware} (Documento 10): la domanda complementare al
Documento 6 --- se il rumore fosse acceso già durante l'ottimizzazione,
l'ansatz convergerebbe altrove? Un argomento fisico (il canale
depolarizzante agisce come trasformazione affine sull'energia, che non
sposta un minimo) prevede di no \emph{prima} di verificarlo; la verifica
--- due metriche, due punti di lavoro, l'intera griglia del Documento 9
--- conferma la previsione ovunque, non solo al punto di riferimento.

\textbf{Readout asimmetrico} (Documento 11): qui la garanzia matematica
del Documento 9 non si estende automaticamente --- il readout asimmetrico
è una trasformazione affine, non più puramente moltiplicativa, e
sommare una costante prima di prendere un modulo non commuta con
l'operazione di modulo. La domanda resta genuinamente aperta fino alla
verifica numerica: $N^*$ risulta invariante, ma è un fatto empirico
specifico di questo correlatore e di questo punto di lavoro, non una
conseguenza necessaria della struttura del problema.

\clearpage
\section{Schema visivo d'insieme}

\begin{figure}[htbp]
\centering
\begin{tikzpicture}[
  box/.style={draw,thick,rounded corners=4pt,align=left,
              font=\footnotesize,inner sep=7pt,
              text width=4.5cm,minimum height=2.6cm},
  extbox/.style={draw,thick,dashed,rounded corners=4pt,align=left,
              font=\footnotesize,inner sep=7pt,
              text width=4.5cm,minimum height=2.6cm},
  lab/.style={font=\scriptsize\itshape,text=gray!70!black,align=center},
  arrow/.style={-{Stealth[length=7pt]},thick},
]

\node[box,draw=blu,fill=blu!5] (n5) at (0,4.8) {
  \textbf{\textcolor{blu}{5. Modello}}\\[2pt]
  2 canali, \texttt{ibm\_torino}\\[2pt]
  \textit{fornisce:} il \texttt{NoiseModel}
};

\node[box,draw=arancio,fill=arancio!7] (n6) at (5.3,4.8) {
  \textbf{\textcolor{arancio!80!black}{6. VQE}}\\[2pt]
  $1{-}F\sim1.5\%$, $2$ CNOT\\[2pt]
  \textit{fornisce:} scala di base
};

\node[box,draw=verdes,fill=verdes!6] (n7) at (10.6,4.8) {
  \textbf{\textcolor{verdes}{7. Dinamica}}\\[2pt]
  $N^*=8$, min. a $N=3$\\[2pt]
  \textit{fornisce:} regola blocchi
};

\node[box,draw=violas,fill=violas!8] (n8) at (10.6,1.4) {
  \textbf{\textcolor{violas}{8. Correlatori}}\\[2pt]
  $N^*=2$, readout a shot\\[2pt]
  \textit{fornisce:} correlatore rif.
};

\node[box,draw=rossos,fill=rossos!6] (n9) at (5.3,1.4) {
  \textbf{\textcolor{rossos}{9. Scan}}\\[2pt]
  $N^*$ vs $\varepsilon$; invariante $p_\text{ro}$\\[2pt]
  \textit{apre:} le estensioni
};

\node[extbox,draw=arancio!70!black,fill=arancio!4] (n10) at (0,-2.0) {
  \textbf{\textcolor{arancio!70!black}{10. Noise-aware}}\\[2pt]
  rioptimizzare: nessun\\
  cambiamento (25/25)
};

\node[extbox,draw=violas,fill=violas!4] (n11) at (5.3,-2.0) {
  \textbf{\textcolor{violas}{11. Readout asim.}}\\[2pt]
  $N^*$ invariato fino a $50\times$\\
  (non garantito dalla teoria)
};

\draw[arrow,blu] (n5) -- (n6);
\draw[arrow,arancio!80!black] (n6) -- (n7);
\draw[arrow,verdes] (n7) -- (n8);
\draw[arrow,violas] (n8) -- (n9);
\draw[arrow,rossos] (n9) -- (n10);
\draw[arrow,rossos] (n9) -- (n11);

\end{tikzpicture}
\caption{Frecce continue: collegamento diretto, stadio per stadio
(Documenti 5--9). Riquadri tratteggiati: le due estensioni facoltative
(Documenti 10, 11), entrambe aperte dallo scan del Documento 9 ma non
collegate fra loro --- due applicazioni indipendenti della stessa
domanda (``$N^*$ si sposta?''), a condizioni diverse.}
\end{figure}

\end{document}
